\section*{Introduction}

% Owner: papers/1_method/approach.md (the framing) and papers/_program/program.md §1 (merged macro
% paper, least squares as the comparison anchor).
%
% LENGTH PASS 2026-08-09. 2025 countable words -> 1771. Target for this section was ~1250 and it is
% NOT met, by 521 words; the shortfall is declared rather than taken out of protected material.
% What binds: six paragraphs were flagged untouchable by the prior triage (the two opening
% paragraphs, the pair on what a least-squares residual can be read for and what modelling the
% correlation takes away, the P2X2 paragraph, the falsifiable ordering, the narrow-novelty claim)
% and they are 704 words on their own. Adding the protected scope clause and the protected
% non-stationary-fluctuation-analysis concession puts the floor near 1000 before any of the
% argument is written. Cutting to 1250 would have required deleting the lineage or the score
% apparatus outright.
% WHAT CAME OUT, all of it verified to survive somewhere else except where marked:
%  - the error-bar conversion table (1.15 -> 7%, 2 -> 41%, 4 -> 95% becomes ~68%), which is stated
%    verbatim at 03_diagnostics.tex:90-92 and belongs there.
%  - "held at fixed spectral density so that shortening the interval raises the noise in each
%    sample and leaves the noise per unit bandwidth unchanged" from the design paragraph. Methods
%    owns it: 06_methods.tex:760-765 defines Current_Noise as a white-noise power spectral density
%    reaching the likelihood only through e_t = Current_Noise/Delta_t. CONSEQUENCE FOR THE BESSEL
%    NOTE below: the clause it calls "the clause at risk" no longer exists here, so the [OPEN]
%    choice it poses is settled by default as "nothing here", and the guard has to land in Methods
%    or in the Discussion's Bessel paragraph. The note is left in place unedited so the decision
%    stays visible.
%  - the ARMA/GLS half of the two-alternatives paragraph, 168 words -> 67. The family, the
%    citation, the empirical finding and the two structural reasons all survive; what went is the
%    free-timescales sentence and the gloss naming N_ch and i as the parameters the fluctuations
%    are richest in, which the sentence before it now carries.
%  - "behaves as though the recording had supplied more observations than it did", which the
%    closing paragraph says in its own words and 05_discussion.tex:137 says with a number.
%  - the Anderson gloss in the lineage paragraph, because the closing paragraph states the same
%    priority claim and states it as the novelty baseline.
%  - "it keeps the exact simulation cheap enough to run thousands of replicates at every design
%    cell": the fact is in the Figure 4 caption (ten thousand recordings per cell) and in
%    06_methods.tex:142.
%  - the delcore gloss ("do not distinguish stochastic channel gating from measurement noise and
%    return biased estimates") and the sweep counts "256 to 504 repeated sweeps". NEITHER IS SAID
%    ELSEWHERE. Both are supporting detail on other people's work, not caveats of ours.
% NOTHING was taken out of a scope caveat, a limitation or a concession. The full list is in the
% pass report.
%
% REWRITTEN 2026-08-04. Predecessor kept verbatim at
% ../archives/01_introduction_SUPERSEDED_20260804.tex (2095 words). This version is ~1700.
% What changed and why, so the trade is not re-litigated:
%
% (a) THE THREAD. The previous version had no single line: it was context, gap, lineage, then four
%     separate defensive digressions before the criterion. This one runs on one, that the memory of
%     the gating process is both what carries the information and what lets a model track the data
%     without earning it. Chosen over four alternatives in a controlled comparison (five drafts on
%     five threads, same fact packet, same length, three critique lenses; 2026-08-04).
% (b) THE OPENING no longer attributes least squares to Hodgkin and Huxley, who fitted by hand, and
%     no longer uses the events-per-sample ratio to separate the single-channel from the macroscopic
%     regime. That separation is the CHANNEL COUNT. The ratio is the interval axis and appears where
%     it belongs.
% (c) THE LINEAGE is compressed from a nine-name walk to the two axes it establishes, with the
%     citations hanging off them. The two members that act in the paper (the non-recursive member,
%     the interval-conditioned one) stay visible.
% (d) THE SCORE PARAGRAPH now says, in the sentence that introduces the identity, that measuring it
%     needs an ensemble from known parameters and is therefore not something a reader runs on one
%     record. That removes an expectation the old closer had to withdraw. It also avoids implying an
%     approximation has a true parameter of its own (the score vanishes at the pseudo-true value,
%     which is not the simulator's).
% (e) GODAMBE 1960 IS OUT. Optimum estimating functions, not the score identity.
% (f) THE ARMA AND FLUCTUATION-ANALYSIS CONCESSIONS MOVED after the criterion. In the predecessor
%     they sat before it and cut the line twice at the point of maximum tension.
% (g) THE CLOSER names the measurement rather than promising it, and raises nothing it then declines.
%     "Locating the reader on the map" was wrong and is gone: the map coordinates come from peak
%     current and baseline variance, while the residual autocorrelation answers a different and
%     narrower question, whether the interval a least-squares fit reports is credible. Whether to use
%     least squares is the experimenter's decision and the paper does not raise it. The demarcation
%     of what is not settled belongs in the Discussion, not here.
%
% STILL OWED, see sections/01_introduction.md:
%  - the "each sample equals the last plus noise" claim needs its derivation (the first difference of
%    a mean-square continuous process vanishes with the interval) or a naive-forecast citation. It is
%    currently an assertion.
%  - covariance fitting cost: "square of the number of samples" (moffatt2007estimation) vs "cube of
%    the record length" (predecessor). Open. The wording below avoids the exponent.
%
% PENDING-BIB: katz1970membrane, katz1972statistical, anderson1973voltage, conti1980conductance,
% sigworth1981covariance and lei2020considering live in biblio_full.bib (124 entries) and are absent
% from biblio.bib (15 entries), which is what elife_paper.tex points at. Repoint \bibliography{biblio}
% to biblio_full, or merge the six keys across. hodgkin1952quantitative, neher1977conductance and
% colquhoun1981stochastic: colquhoun1981stochastic and neher1977conductance are in biblio_full;
% hodgkin1952quantitative is in NEITHER and must be added.

A channel opens and shuts at random, and how long it stays in each conformation is the elementary
random quantity of gating. What an experiment sees of it depends on how many channels contribute.
With one channel in the patch the durations are read off the record and the rate constants come from
their distributions \citep{neher1977conductance,colquhoun1981stochastic}. In the currents that carry
an action potential so many channels contribute that the durations average away, the current follows
the deterministic relaxation of the Hodgkin and Huxley description \citep{hodgkin1952quantitative},
and rate constants are obtained by fitting that relaxation, in current practice by least squares.

Most macroscopic patch-clamp recordings are in neither place. A few hundred to a few thousand
channels contribute, individual openings cannot be resolved, and the averaging is incomplete. The
durations are still there, and what they leave behind is a fluctuation around the mean current that
carries the memory of the gating process. Samples taken closer together than the channel's
relaxation time are correlated, so a sample taken faster than that time partly repeats what the one
before it already said.

That fluctuation carries what the mean cannot. Channel number and unitary current enter the mean
current as a product, so no fit to the mean separates them. The standard analysis discards the
fluctuation anyway: the four approaches surveyed by \citet{clerx2019four} are deterministic fits with
a root-mean-square objective, and the IonBench benchmark states that it does not cover optimisation
approaches for stochastic models \citep{owen2025ionbench}. Discarding it usually leaves the point
estimate standing. What it damages is the count of independent residuals, the number that turns a
residual sum of squares into a confidence interval. \citet{delcore2025parameter} state the same cost in the field's own terms, that standard
methods built for deterministic models do not distinguish stochastic channel gating from
measurement noise and return biased estimates.
% What the classical arm of this paper is, said plainly and once: least squares on the deterministic
% mean current with the unitary current held fixed. Repeated in Methods; do not let it drift.

A Markov description models the dependence rather than assuming it away. A channel with a small
number of conformational states, connected by rate constants that hold over the whole record,
generates the mean current and the covariance of the fluctuations from the same constants, so the
temporal structure of the noise becomes a second measurement of the same parameters. Those constants
also mean something outside the record: they are properties of the protein, comparable across
preparations and laboratories, and they connect what is measured to a mechanism. A correlation
function fitted freely to the residuals delivers neither of those.

Kinetic information has been read from that correlation since the 1970s, first through stationary
power spectra \citep{katz1970membrane,katz1972statistical,anderson1973voltage} and then through the
exact non-stationary two-time covariance of a population of identical independent channels,
estimated from ensembles of $256$ to $504$ repeated sweeps
\citep{conti1980conductance,sigworth1981covariance}. The whole-record likelihoods that followed
differ along two axes. The first is how much of the gating stochasticity the likelihood models:
nothing, then the variance at each sample \citep{milescu2005maximum}, then its correlation in time,
whether by carrying the full two-time covariance at a cost that keeps it to a hundred or two selected
points \citep{celentano2004use}, or by propagating the occupancy recursively and conditioning it on
each new datum at a cost linear in the record
\citep{moffatt2007estimation,stepanyuk2011efficient,munch2022bayesian}. The second axis is the
observable: an acquisition system reports an average over each interval rather than an instantaneous
sample, and the most recent member conditions that interval average on the channel states at both of
its ends \citep{moffatt2025bayesian}.

None of these has displaced least squares. Cost is part of it, and \citet{delcore2025parameter} make
that case. There is also something these methods ask an analyst to give up. A least-squares fit
predicts the current at every point from the parameters alone, so the residual can be read directly,
and read twice. Its amplitude reports the quality of the fit against what the gating variance and the
baseline noise should deliver, and its correlation in time reports what the fit leaves out. Both
readings appeal to knowledge from outside the fit, which is what makes them checks.

Modelling the correlation is what takes those readings away. A recursive likelihood predicts each
interval from the parameters and from the data already observed, so the previous sample enters the
prediction of the next one and the predicted trace follows the record by construction. It follows the
record when the scheme is right, when the scheme is wrong, and when the code has a mistake in it. The
limiting case of the argument carries no kinetics at all and states that each sample equals the last
one plus noise: where the interval is short against the relaxation time and the instrumental noise is
small, it predicts almost perfectly and says nothing about the channel. % NEEDS SUPPORT, see header.
Such a likelihood also standardises its residual against its own predicted variance, so the amplitude
reading becomes a statement of internal consistency, and the correlation reading is spent because the
recursion subtracts at each step what it has just seen. The same blindness covers two questions that
an experimenter cannot separate and has no reason to: whether the approximation is adequate, and
whether the code that computes it is correct.

That is not a hypothetical. Nine kinetic schemes for the P2X$_2$ receptor were ranked by Bayesian
evidence twice over the same outside-out patch recordings, in work of our own, and the two rankings
were not the same \citep{moffatt2025bayesian}. What separated them was the recursion. Both
likelihoods averaged the current over the acquisition interval and both carried the same gating
variance; one conditioned each interval on the data already observed and the other did not. That the
ranking moved with the approximation was reported. Which of the two orderings to believe was not,
because it turns on how far each likelihood's reported information sits from the information it
delivers, and that had never been measured for either.
% ADDED 2026-08-07, 116 words. Owner: the desk-rejection pass recorded at 00_abstract.tex note 1f.
% WHY HERE. The paragraph above ends on the two questions an experimenter cannot separate and then
% went straight to the classical apparatus, so "this" in "The apparatus that answers this" had an
% abstract antecedent. This paragraph is that same sentence with a date and a receptor, and it hands
% the reader to the apparatus with a concrete thing for it to answer. It rides the section's thread
% rather than cutting it: the thread is that the memory of the gating process is both what carries
% the information and what lets a model track the data without earning it, and this is the case where
% the unearned tracking reached a mechanism.
% NO NEW CONCEPTS. Recursion, the acquisition interval, conditioning on data already observed and the
% gating variance are all defined in the lineage paragraph above, which already cites this same
% paper for the interval-conditioned member.
% THE AXIS DISCIPLINE IS LOAD-BEARING, see 00_abstract.tex note 1f(i). "What separated them was the
% recursion" is exact: the two runs behind Figs 1d and 1e of that paper shared averaging over the
% acquisition interval and the same gating variance term and differed in the recursion alone. Never
% write "MacroIR against MacroINR" or "the boundary-conditioned member" here. The object that ran
% there is not the member specified in Theory, so naming it as IR would be false, and the axis
% wording is both true and the axis this paper measures.
% "WERE NOT THE SAME", NOT "PUT DIFFERENT MECHANISMS FIRST", which is the stronger sentence: that
% paper numbers nine schemes I-IX in its main text and eleven I-XI in Table S1 with no published
% mapping, so the identity of the top-ranked scheme under each likelihood cannot be checked from the
% record. Upgrade only on a first-hand check.
% WHAT THIS PARAGRAPH MUST NOT DO is adjudicate. It says the ordering moved and that nothing in that
% analysis could say which to believe. The matching declining sentence is at 05_discussion.tex, in
% the paragraph that carries the evidence bridge; the two must stay consistent.

The apparatus that answers this is classical. For a correctly specified likelihood the score, the
gradient of the log-likelihood with respect to the parameters, has mean zero at the parameters that
generated the data, and its covariance across independent recordings equals the Fisher information
the likelihood reports \citep{huber1967behavior,white1982maximum}. An approximation satisfies neither
statement at those parameters, and the two departures are measurable quantities rather than tests to
be passed. Goodness of fit does not answer the question: an approximation can reproduce the mean
current closely and still report an information matrix wrong by more than an order of magnitude.
Comparing two approximations against each other does not answer it either, because both may be
wrong. Measuring the departures takes an ensemble of recordings drawn from known parameters, which is
why this is a simulation study and not something a reader can run on a single record. For this class
of model it is possible at all because the forward process can be simulated exactly while its
likelihood cannot, so the simulator supplies ground truth for what the likelihood approximates.

Two alternatives to a mechanistic likelihood are worth placing here, since both are in use. The
standard repair for correlated residuals is generalized least squares with autoregressive or
moving-average errors, which whitens the residual sequence with a correlation function estimated
from the residuals themselves, and \citet{lei2020considering} ran that on ion-channel data and
report that it widens the posterior without predicting better. Gating puts that family out of
reach: such an error model is stationary by construction while the gating covariance tracks the
mean current and restarts at every concentration or voltage jump, and it carries no channel number
and no unitary current. Its timescales are free where the Markov model ties them to the eigenvalues
of the rate matrix that already fixed the mean, so they absorb the information the mechanism would
have contributed. The other alternative is
non-stationary fluctuation analysis, the one fluctuation-based method in routine use, which regresses
the variance of the current against its mean and returns the unitary current, the channel number and
the peak open probability \citep{stepanyuk2014maximum}. This paper does not compete with it over how
much amplitude information a macroscopic recording holds, never benchmarks against it, and no
comparison with the variance-mean regression should be read into any result below.

Because every member is an approximation, the measurement returns a domain rather than a verdict: a
map of where a likelihood reports honestly, bounded by level sets of a continuous diagnostic rather
than by thresholds. Position on it is set by the number of channels, the acquisition interval in
units of the channel's relaxation time, and the instrumental noise in units of the gating noise. The
measurements are made on a two-state scheme at an open probability of one half, driven by a single
concentration jump, compared at the level of the likelihood rather than of a fitted experiment. The
small model isolates the error the approximation itself commits, free of confounding from a
misspecified kinetic scheme, and it is the most favourable case for a Gaussian treatment of the
occupancy, therefore a lower bound on the error for anything richer. Open probability, stationary
protocols and larger schemes are left out for reasons that follow from the same design, and are
stated with the results.
% [BESSEL-BOUND] OPEN 2026-08-06. Canonical note at 02_theory.tex, under the Bessel paragraph.
% READ WITH THE LENGTH-PASS NOTE IN THE HEADER: the clause this note guards was removed on
% 2026-08-09 as specification detail Methods owns, so the [OPEN] question below is settled here by
% default and the guard has to land at 06_methods.tex or in the Discussion's Bessel paragraph.
% Note kept verbatim so the decision is not lost.
% THE CLAUSE AT RISK is "held at fixed spectral density so that shortening the interval raises the
% noise in each sample and leaves the noise per unit bandwidth unchanged". That is exactly right
% for the simulator and it has a physical ceiling on a real rig, which the sentence does not say.
% A real amplifier's noise is flat only up to its own anti-alias filter: once the window falls
% below the filter's response time the per-sample variance stops rising as 1/Delta and the excess
% appears as correlation between neighbouring samples instead. So the interval axis of this study
% is realisable on a bench down to about Delta = a few / f_c and no further, unless the filter is
% opened with it, which costs noise.
% WHY IT IS WORTH A CLAUSE RATHER THAN A PARAGRAPH. The same dimensionless group, f_c*Delta,
% governs both this and the observable's kernel (canonical note, items i and ii), so the two
% departures are one departure seen from two sides, and both vanish at the same rate. Nothing in
% the paper's measurements changes: the simulator's noise IS white by construction and the sweep
% is over a model, which this paragraph already says. What is missing is that the shortest
% intervals of the sweep are the ones a bench cannot deliver at fixed spectral density, and that
% is the same corner where the boundary-conditioned member's advantage is largest.
% [OPEN] Either a subordinate clause here ("... unchanged, which a real amplifier delivers down
% to its anti-alias filter's response time and no further") or nothing here and one sentence in
% the Discussion's Bessel paragraph. Do NOT put it in both; this paragraph is already the densest
% in the section.

Two things come out of the map, and they answer different readers. The first is where each
approximation misreports what the data determine, including the most complete one, since a method
whose limits are unmapped has been endorsed rather than characterised. The second is the region in
which a classical fit reports an interval that matches the spread it delivers, and it carries an
ordering that is falsifiable and does not depend on where the level sets are placed: there is no
region of the measured plane in which the fluctuations still carry information about the unitary
current or the channel number and the classical interval is simultaneously trustworthy, the two
boundaries being separated by one to three decades of instrumental noise everywhere they were
measured. % src: projects/eLife_2025/figures/paper_both/figure_6.html (crossing table, fitted slopes)

Against that history, what is new here is narrow and is worth stating narrowly. A joint estimate of
rates and amplitudes from fluctuation data, carrying standard errors, is \citet{anderson1973voltage}.
What this paper adds, to a non-stationary protocol read from a single record, is the third element of
the combination: an interval whose calibration has been verified rather than assumed. The thousands of
replicate recordings used below certify the method; they are not what applying it requires. Separately
from the map, there is a reading an experimenter can take from a recording they already have. The
integrated autocorrelation of the standardized least-squares residual counts how many intervals the
fit effectively has for every interval it believes it has, and it costs nothing, since the fit was
going to be made anyway. Two design variables that look as though they should matter, the channel
count and the instrumental noise, enter the answer through that one statistic and through nothing
else. The acquisition interval does not, and the reader knows it, because they chose it when they set
the sampling rate, as they did the analogue filter, which contributes to the same statistic on a
record sampled faster than its own bandwidth. Those two numbers together say what a classical fit is
doing to the width of the interval it reports.
% src: projects/eLife_2025/figures/paper_both/figure_6B.Rmd (header, measured points 1 to 5): tau_int
% in normalized Sokal/Kish form, null 1 for a white residual; at matched tau_int the least-squares
% size agrees to three digits over four decades of N_ch (1.127 / 1.135 / 1.123 / 1.145 at ~1.04), while
% at matched tau_int two sampling intervals differ by a factor 5.7 (8.9 at interval 0.01 against
% 1.56-1.63 at 0.5). The error-bar scale is the square root of the distortion, which is why 4 and not 3
% is the factor-two mark.
% LENGTH PASS 2026-08-09: the conversion table this src note documents ("in units a reader can act
% on: a distortion of 1.15 is a 7% error on a confidence interval, 2 is 41%, and 4 turns a nominal
% 95% interval into about 68%") was cut from the sentence above because 03_diagnostics.tex:90-92
% states it verbatim, in the paragraph that defines the distortion it converts. The src note stays
% here because the sentence it also supports, on what the two numbers say about the width of the
% reported interval, stays.
% CORRECTED 2026-08-04: an earlier draft of this paragraph said the residual autocorrelation alone
% locates a recording, and quoted a factor of two at three. Both were wrong, and figure_6B.Rmd had
% already measured the right version.
% [BESSEL-BOUND] CLOSED HERE 2026-08-07 by the fifteen-word clause the note below asked for
% ("as they did the analogue filter, which contributes to the same statistic on a record sampled
% faster than its own bandwidth"). The other two sites, 03_diagnostics.tex beside the
% residual-whiteness check and the scope clause in the design paragraph above, are still open, and
% the canonical note with the bound and the script stays at 02_theory.tex under the Bessel
% paragraph. What was bought: the paragraph no longer asserts that two design variables enter the
% answer and nothing else does, on a record where a third one demonstrably enters, and the repair
% rides the sentence's own logic rather than contradicting it, because the filter is a number the
% reader chose exactly as they chose the sampling rate. What was NOT bought and must not be read
% in: no size for the contribution is stated here, deliberately, because stating it commits to the
% expansion, which belongs in Theory. Original note follows.
% [BESSEL-BOUND] OPEN 2026-08-06. Canonical note at 02_theory.tex; the argument for why this
% particular sentence needs guarding is at 03_diagnostics.tex, beside the residual-whiteness
% check.
% THIS PARAGRAPH MAKES THE PAPER'S ONE OFFER TO A READER WITH DATA IN HAND: "the integrated
% autocorrelation of the standardized least-squares residual ... costs nothing, since the fit was
% going to be made anyway". On a real recording that statistic also picks up the analogue filter,
% which contributes a lag-one residual correlation of 0.075/(f_c*Delta) on its own, 8.6% where
% the acquisition window equals the filter's inverse cutoff. The offer is sound where the record
% is sampled at or below the filter's bandwidth and misleading where it is not, and the reader
% cannot tell the two apart from the number alone: a filter and a badly correlated fit both push
% tau_int up.
% ALSO EXPOSED, and it is the same fix: "Two design variables that look as though they should
% matter, the channel count and the instrumental noise, enter the answer through that one
% statistic and through nothing else." On an oversampled record a THIRD variable enters it, the
% filter setting, and that one the reader also chose, like the sampling rate. The sentence's own
% logic therefore extends to it rather than being broken by it, which is the cheapest repair:
% name the filter alongside the sampling rate in the clause that already says the reader knows
% the acquisition interval because they set it.
% COST: about fifteen words. This is the highest-value fifteen words the Bessel scoping buys
% anywhere in the manuscript, because it is the only place a reader is told to act.
